% ============================================================================
% Core Feature: Approval Workflow
% ============================================================================
\section{Approval Workflow}

\begin{ugModule}{Overview}

  \textbf{Purpose:}
  Route documents through a multi-level, position-based
  approval chain with full audit trail and rejection handling.

  \textbf{Who Can Access:}
  \ugRole{Admin} \ugRole{Manager} \ugRole{Reviewer}
\end{ugModule}

% ────────────────────────────────────────────────────────────────
\subsection{How Approval Routing Works}

Unlike a simple status-based workflow, InsightDoc uses a
\textbf{position-based multi-level approval chain}. Each
document is routed through up to 6 approval levels, where
each level is assigned to a specific unit division position.

The approval chain is determined by the \ugField{Approval
Purpose} selected during document creation (see
Document Management, Tab~1). Each Approval Purpose defines:

\begin{itemize}[leftmargin=18pt]
  \item Which \textbf{Unit} it belongs to
  \item Up to \textbf{6 approval levels}, each mapped to a
        Unit Division Level (position/role within the unit)
  \item The \textbf{sequential order} of approvals
\end{itemize}

\begin{ugNote}
  Approval Purposes are configured by the \ugRole{Admin} in
  \ugMenu{Settings > Approval Configuration > Approval
  Purpose}. Each purpose defines a complete signing chain
  specific to a unit and document type.
\end{ugNote}

% ────────────────────────────────────────────────────────────────
\subsection{Approval Levels}

Each document can have up to 6 approval levels:

\begin{tabularx}{\textwidth}{@{} c l X @{}}
  \toprule
  \textbf{Level} & \textbf{Position}
    & \textbf{Description} \\
  \midrule
  1 & Level 1 Approver
    & First approval in the chain (e.g., Section Head) \\
  2 & Level 2 Approver
    & Second approval (e.g., Department Head) \\
  3 & Level 3 Approver
    & Third approval (e.g., Division Head) \\
  4 & Level 4 Approver
    & Fourth approval (e.g., Director) \\
  5 & Level 5 Approver
    & Fifth approval (e.g., VP) \\
  6 & Level 6 Approver
    & Final approval (e.g., President Director) \\
  \bottomrule
\end{tabularx}

\vspace{6pt}

Not all levels are required. If an Approval Purpose only
defines Level~1 through Level~3, levels 4--6 are skipped
automatically.

% ────────────────────────────────────────────────────────────────
\subsection{Approval Status Per Level}

Each approval level has its own status, independent of the
overall document status:

\begin{tabularx}{\textwidth}{@{} l l X @{}}
  \toprule
  \textbf{Status} & \textbf{Code}
    & \textbf{Description} \\
  \midrule
  Waiting Previous
    & \texttt{waiting-previous}
    & This level is blocked until all previous
      levels are approved \\
  Waiting Approval
    & \texttt{waiting-approval}
    & This level is the current active approver;
      action is required \\
  Approved
    & \texttt{approved}
    & This level has been approved; recorded with
      approver name and timestamp \\
  Rejected
    & \texttt{rejected}
    & This level has been rejected; recorded with
      rejector name, timestamp, and reason \\
  Skipped
    & \texttt{skipped}
    & This level was not configured for this
      document's Approval Purpose \\
  \bottomrule
\end{tabularx}

% ────────────────────────────────────────────────────────────────
\subsection{Approval Flow}

The approval process follows this sequential logic:

\begin{enumerate}[leftmargin=18pt]
  \item When a document is created, the system generates
        \texttt{DocumentApproval} records for each configured
        level.
  \item Level~1 starts with
        \texttt{waiting-approval} status.
  \item Levels 2+ start with
        \texttt{waiting-previous} status.
  \item Unconfigured levels are set to
        \texttt{skipped}.
  \item When Level~$n$ approves, Level~$n+1$ transitions
        from \texttt{waiting-previous} to
        \texttt{waiting-approval}.
  \item When all configured levels approve, the document
        status transitions to
        \ugStatus{ugInfo}{Approved}.
  \item If \textbf{any} level rejects, the entire document
        is rejected --- the document is marked as
        \ugStatus{ugDanger}{Rejected} with the rejection
        reason, rejector identity, and the level that rejected.
\end{enumerate}

% ────────────────────────────────────────────────────────────────
\subsection{Status Progress Track}

On the Document List page, the \ugField{Status} column displays
a visual \textbf{multi-step progress track} showing all
approval levels and their current status:

\begin{itemize}[leftmargin=18pt]
  \item Each step represents one approval level
  \item Approved levels show a green checkmark
  \item The current waiting level is highlighted
  \item Rejected levels show a red indicator with the
        rejection reason visible on hover
  \item Skipped levels are grayed out
\end{itemize}

This gives users an immediate visual overview of where a
document stands in its approval chain.

% ────────────────────────────────────────────────────────────────
\subsection{Approval Data Recorded}

Each \texttt{DocumentApproval} record stores the following:

\begin{tabularx}{\textwidth}{@{} l X @{}}
  \toprule
  \textbf{Field} & \textbf{Description} \\
  \midrule
  Document
    & Reference to the document being approved \\
  Level
    & The approval level (1--6) \\
  Unit Division
    & The organizational position assigned as
      approver \\
  Approval Status
    & Current status of this level \\
  User
    & The specific user assigned to approve \\
  Approved By
    & Name of the person who approved \\
  Approved Date
    & Timestamp of approval \\
  Rejected By
    & Name of the person who rejected \\
  Rejected Date
    & Timestamp of rejection \\
  Rejected User
    & Reference to the rejecting user's account \\
  \bottomrule
\end{tabularx}

% ────────────────────────────────────────────────────────────────
\subsection{Rejection Handling}

When a document is rejected at any level:

\begin{enumerate}[leftmargin=18pt]
  \item The rejection is recorded on both the
        \texttt{DocumentApproval} record (level-specific) and
        the \texttt{Document} entity itself.
  \item The document stores:
    \begin{itemize}[leftmargin=14pt]
      \item \texttt{IsRejected = true}
      \item \texttt{RejectedLevelId} --- which level rejected
      \item \texttt{RejectedBy} --- who rejected
      \item \texttt{RejectedReason} --- mandatory comment
      \item \texttt{RejectedDate} --- when
      \item \texttt{RejectedUserId} --- rejector's user ID
    \end{itemize}
  \item The document's \ugField{Document Status} transitions to
        \ugStatus{ugDanger}{Rejected}.
  \item The document owner is notified and can revise the
        document for resubmission.
\end{enumerate}

\begin{ugImportant}
  Rejection reasons are \textbf{mandatory} and permanently
  recorded. They are visible on the document's status progress
  track (on hover) and in the document detail page. Write
  specific, actionable feedback to help the document owner
  address the issues effectively.
\end{ugImportant}

% ────────────────────────────────────────────────────────────────
\subsection{Step-by-Step: Approving a Document}

\begin{ugSteps}
  \item Navigate to \ugMenu{Approval Center}.
  \item The list shows documents where you are the current
        active approver (your level has
        \texttt{waiting-approval} status).
  \item Click a document to view its detail and preview.
  \item Review the document content via the in-browser PDF
        preview.
  \item Click \ugButton{Approve} to approve at your level.
  \item Confirm the action in the modal dialog.
  \item If you are the final approver, the document transitions
        to \ugStatus{ugInfo}{Approved}. Otherwise, the next
        level is activated.
\end{ugSteps}

\subsection{Step-by-Step: Rejecting a Document}

\begin{ugSteps}
  \item Navigate to \ugMenu{Approval Center}.
  \item Click the document to review.
  \item Click \ugButton{Reject}.
  \item Enter a \textbf{mandatory rejection reason} in the
        comment field.
  \item Confirm the rejection.
  \item The document is marked as Rejected with your identity,
        the level, and the reason recorded.
\end{ugSteps}

% ────────────────────────────────────────────────────────────────
\subsection{Activity Tracking}

All approval and rejection actions are recorded as
\texttt{DocumentActivity} entries with the following action
codes:

\begin{tabularx}{\textwidth}{@{} l X @{}}
  \toprule
  \textbf{Action Code} & \textbf{Description} \\
  \midrule
  \texttt{approve-document}
    & Document was approved at a specific level \\
  \texttt{reject-document}
    & Document was rejected at a specific level \\
  \texttt{upload-document}
    & A new document was created or revised \\
  \bottomrule
\end{tabularx}

Each activity record includes: user identity, browser, OS,
device, IP address, timestamp, and a unique request ID for
traceability.

% ────────────────────────────────────────────────────────────────
\subsection{Business Rules}

\begin{itemize}[leftmargin=18pt]
  \item Only the user assigned to the current active level
        can approve or reject.
  \item Approval is \textbf{strictly sequential} --- Level~2
        cannot act until Level~1 approves.
  \item A rejection at \textbf{any} level rejects the entire
        document.
  \item Rejected documents can be revised and resubmitted,
        restarting the approval chain from Level~1.
  \item Skipped levels do not block the approval flow.
  \item The Approval Purpose configuration determines which
        levels are active; this cannot be changed after
        document creation.
  \item All approval actions are immutable audit records.
\end{itemize}

\begin{ugTip}
  As a reviewer, check your Approval Center daily. Documents
  waiting for your approval block the entire chain for all
  subsequent approvers.
\end{ugTip}

\newpage
